\section{Introduction}
Les systèmes robotiques auto-équilibrés constituent un domaine d'application particulièrement intéressant de la robotique moderne, car ils combinent des problématiques de mécanique, d'électronique, d'informatique embarquée et de théorie du contrôle. Parmi ces systèmes, le robot à deux roues auto-balançant représente un excellent cas d'étude en raison de son instabilité intrinsèque et des défis techniques qu'il implique.

Le principe de fonctionnement d'un tel robot repose sur la capacité à maintenir en permanence son équilibre malgré sa configuration naturellement instable, comparable à celle d'un pendule inversé. Pour y parvenir, le système doit mesurer son inclinaison à l'aide de capteurs, traiter les données en temps réel et ajuster la vitesse des moteurs afin de corriger les écarts observés. Cette boucle de rétroaction nécessite une intégration étroite entre les composants matériels et les algorithmes de commande.

L'objectif de ce travail de Bachelor est de concevoir, fabriquer et valider un prototype fonctionnel de robot à deux roues auto-balançant. Le projet couvre l'ensemble des étapes de développement, depuis la conception mécanique et le choix des composants électroniques jusqu'à l'implémentation des algorithmes de contrôle et la réalisation des essais expérimentaux.

Dans un premier temps, les principes théoriques liés au pendule inversé et aux systèmes asservis sont présentés. Ensuite, l'architecture matérielle du robot est détaillée, notamment les moteurs, les capteurs inertiels et le microcontrôleur utilisés. Une attention particulière est portée au développement de l'algorithme de stabilisation permettant d'assurer le maintien de l'équilibre. Enfin, des tests expérimentaux sont réalisés afin d'évaluer les performances du système et d'identifier les éventuelles améliorations à apporter.

Au-delà de l'aspect pédagogique, ce projet illustre des problématiques largement rencontrées dans le domaine de la robotique mobile et des véhicules autonomes. Les compétences acquises au cours de ce travail constituent ainsi une base solide pour la compréhension et le développement de systèmes robotiques plus complexes.

\section{Contexte}

La robotique mobile connaît un développement important, notamment dans les domaines de l'automatisation, des transports intelligents et des systèmes embarqués. Parmi les différentes plateformes robotiques, les systèmes auto-équilibrés occupent une place particulière en raison de leur dynamique instable et des défis de contrôle qu'ils impliquent.

Le robot à deux roues auto-balançant est un exemple classique de système dit \textit{non linéaire instable}, souvent modélisé comme un pendule inversé. Ce type de système est largement utilisé dans l'enseignement et la recherche en automatique, car il permet d'illustrer concrètement des concepts tels que la rétroaction, la stabilisation et la commande en temps réel.

Dans un contexte industriel et académique, ce type de robot constitue également une base technologique proche de certaines applications réelles, telles que les véhicules de mobilité personnelle auto-équilibrés ou certaines plateformes robotiques autonomes. Il représente donc un excellent compromis entre complexité théorique et réalisation pratique.

Ce travail s'inscrit dans cette perspective en proposant la conception et la réalisation d'un prototype fonctionnel, permettant d'explorer les aspects théoriques et pratiques liés à la stabilisation d'un système instable. L'ensemble du projet mobilise des compétences en mécanique, électronique et informatique embarquée, avec une attention particulière portée à la mise en oeuvre d'un système de contrôle robuste et réactif.
\section{Citations et bibliographie}

\cite{eth}